\documentclass[preview]{standalone}

\usepackage{amsmath}
\usepackage{amssymb}
\usepackage{stellar}
\usepackage{definitions}

\begin{document}

\id{ring-homomorphisms-exercises}
\genpage

\section{Exercises}

\begin{snippetexercise}{quadratic-integer-ring-matrix-representation-exercise}{A matrix representation of a quadratic integer ring}
    Let \(d\in\integers\) be square-free and nonsquare, and let
    \[
        A=\integers[\sqrt d]
        =\{a+b\sqrt d\suchthat a,b\in\integers\}.
    \]
    Define
    \[
        \varphi\colon A\fromto\matrices_2(\integers),
        \qquad
        \varphi(a+b\sqrt d)
        =
        \begin{pmatrix}
            a&bd\\
            b&a
        \end{pmatrix}.
    \]
    Show that \(\varphi\) is an injective unital
    \ringhomomorphism.
\end{snippetexercise}

\begin{snippetsolution}{quadratic-integer-ring-matrix-representation-solution}{A matrix representation of a quadratic integer ring}
    The nonsquare assumption makes the representation
    \(a+b\sqrt d\) unique, so \(\varphi\) is well-defined. It sends
    \(0\) to the zero matrix and \(1\) to \(I_2\). If
    \[
        \alpha=a_1+b_1\sqrt d,
        \qquad
        \beta=a_2+b_2\sqrt d,
    \]
    then direct calculation gives
    \[
        \varphi(\alpha-\beta)=\varphi(\alpha)-\varphi(\beta)
    \]
    and
    \begin{align*}
        \varphi(\alpha)\varphi(\beta)
        &=
        \begin{pmatrix}
            a_1a_2+b_1b_2d&(a_1b_2+a_2b_1)d\\
            a_1b_2+a_2b_1&a_1a_2+b_1b_2d
        \end{pmatrix}\\
        &=\varphi(\alpha\beta).
    \end{align*}
    Thus \(\varphi\) is a unital ring homomorphism.

    If \(\varphi(a+b\sqrt d)=0\), its diagonal and lower-left entries
    give \(a=b=0\). Hence
    \[
        \ringker\varphi=\{0\},
    \]
    and \(\varphi\) is injective. Consequently,
    \[
        \integers[\sqrt d]
        \cong
        \left\{
            \begin{pmatrix}a&bd\\b&a\end{pmatrix}
            \suchthat a,b\in\integers
        \right\}
        \subringleq\matrices_2(\integers).
    \]
\end{snippetsolution}

\begin{snippetexample}{gaussian-integers-real-matrix-representation-example}{Gaussian integers and complex numbers}
    For \(d=-1\),
    \[
        \integers[i]
        \cong
        \left\{
            \begin{pmatrix}a&-b\\b&a\end{pmatrix}
            \suchthat a,b\in\integers
        \right\}.
    \]
    Replacing integer coefficients with real coefficients gives
    \[
        \complexnumbers
        \cong
        \left\{
            \begin{pmatrix}x&-y\\y&x\end{pmatrix}
            \suchthat x,y\in\realnumbers
        \right\}
        \subringleq\matrices_2(\realnumbers).
    \]
\end{snippetexample}

\end{document}
